\documentclass[12pt]{article}

\usepackage[a4paper,margin=2.5cm]{geometry}
\usepackage{graphicx}
\usepackage{booktabs}
\usepackage{tcolorbox}
\usepackage{xcolor}
\usepackage{hyperref}
\usepackage{listings}

\definecolor{primary}{HTML}{1F4E79}
\definecolor{softblue}{HTML}{E8F1FA}

\title{\textbf{\textcolor{primary}{JustBeat-DAW}}\\
\large Documento Maestro de Ingeniería\\
8-bit Digital Audio Workstation}
\author{Equipo de Desarrollo}
\date{\today}

\begin{document}

\maketitle
\tableofcontents
\newpage

%====================================================
\section{Visión General del Sistema}

JustBeat-DAW es una estación de trabajo de audio digital enfocada en la producción musical 8-bit (chiptune), diseñada bajo principios formales de ingeniería de software.

El sistema combina:
\begin{itemize}
    \item Síntesis digital retro.
    \item Arquitectura limpia.
    \item Extensibilidad mediante plugins.
    \item Soporte MIDI profesional.
\end{itemize}

%====================================================
\section{Arquitectura Global}

\subsection{Modelo Arquitectónico}

Se adopta una combinación de:

\begin{itemize}
    \item Clean Architecture
    \item Arquitectura en Capas
    \item Arquitectura basada en eventos
    \item Audio Graph orientado a nodos
\end{itemize}

\subsection{Capas del Sistema}

\begin{tcolorbox}[colback=softblue,colframe=primary,title=Capas Principales]

\textbf{Domain Layer}  
Entidades musicales, eventos MIDI, contratos de plugins.

\textbf{Application Layer}  
Casos de uso: reproducción, grabación, exportación, sincronización.

\textbf{Infrastructure Layer}  
Motor de audio, adaptadores MIDI, sistema de archivos, cargador de plugins.

\textbf{Presentation Layer}  
Interfaz gráfica (PyQt6) y controladores.

\end{tcolorbox}

Regla: Las dependencias apuntan siempre hacia el dominio.

%====================================================
\section{Modelo de Dominio Completo}

\subsection{Entidades Principales}

\begin{itemize}
    \item Project
    \item Track
    \item Pattern
    \item Step
    \item Note
    \item MIDIEvent
    \item Oscillator
    \item Envelope
    \item AudioBuffer
    \item PluginDescriptor
    \item AudioNode
    \item AudioBus
\end{itemize}

%====================================================
\section{Audio Engine Profesional}

\subsection{Componentes Internos}

\begin{itemize}
    \item Clock Engine
    \item Step Scheduler
    \item Oscillator Bank
    \item Envelope Processor
    \item Audio Mixer
    \item Master Bus
\end{itemize}

\subsection{Audio Graph}

El motor se basa en un grafo dirigido:

\begin{tcolorbox}[colback=softblue,colframe=primary,title=Audio Flow]

MIDI Input → Instrument Plugin → Effect Chain → Track Bus → Master Bus → Audio Output

\end{tcolorbox}

Cada nodo implementa:

\begin{lstlisting}
class AudioNode:
    def process(self, buffer):
        pass
\end{lstlisting}

%====================================================
\section{Sistema MIDI Completo}

\subsection{Funcionalidades}

\begin{itemize}
    \item Detección de dispositivos.
    \item Grabación en tiempo real.
    \item Cuantización automática.
    \item MIDI Clock Sync.
    \item Exportación .mid.
\end{itemize}

\subsection{Arquitectura MIDI}

Se implementa un Adapter Pattern:

\begin{lstlisting}
class IMIDIAdapter:
    def connect(self, device_id): pass
    def listen(self): pass
\end{lstlisting}

Infraestructura recomendada:
\begin{itemize}
    \item mido
    \item python-rtmidi
\end{itemize}

%====================================================
\section{Sistema de Plugins Avanzado}

\subsection{Plugin Host}

El sistema implementa un Plugin Host desacoplado.

\begin{itemize}
    \item Carga dinámica con importlib.
    \item Registro automático.
    \item Validación por interfaz.
\end{itemize}

\subsection{Contrato Base}

\begin{lstlisting}
class IPlugin:
    def initialize(self): pass
    def process(self, buffer): pass
    def shutdown(self): pass
\end{lstlisting}

\subsection{Tipos}

\begin{itemize}
    \item InstrumentPlugin
    \item EffectPlugin
    \item UtilityPlugin
\end{itemize}

\subsection{Aislamiento}

Opcional:
\begin{itemize}
    \item multiprocessing para sandbox lógico.
\end{itemize}

%====================================================
\section{Sistema de Routing y Buses}

\subsection{Track Bus}

Cada pista tiene:
\begin{itemize}
    \item Cadena de efectos.
    \item Control de volumen.
    \item Control de paneo.
\end{itemize}

\subsection{Master Bus}

\begin{itemize}
    \item Compresión global futura.
    \item Limitador.
    \item Exportación final.
\end{itemize}

%====================================================
\section{Principios SOLID Aplicados}

\subsection{S}
AudioEngine no gestiona GUI.

\subsection{O}
Nuevos osciladores se agregan sin modificar el núcleo.

\subsection{L}
Todo Oscillator respeta interfaz base.

\subsection{I}
Interfaces pequeñas:
\begin{itemize}
    \item IAudioNode
    \item IMIDIAdapter
    \item IPlugin
\end{itemize}

\subsection{D}
Application depende de abstracciones.

%====================================================
\section{Patrones de Diseño Implementados}

\begin{itemize}
    \item Factory → Creación de osciladores.
    \item Strategy → Algoritmos de síntesis.
    \item Observer → Eventos MIDI.
    \item Adapter → Integración MIDI.
    \item Repository → Persistencia.
    \item Command → Undo/Redo.
    \item Dependency Injection → Desacoplamiento.
\end{itemize}

%====================================================
\section{Sistema Undo/Redo}

Se implementa Command Pattern formal:

\begin{lstlisting}
class ICommand:
    def execute(self): pass
    def undo(self): pass
\end{lstlisting}

Acciones soportadas:
\begin{itemize}
    \item Agregar nota.
    \item Eliminar nota.
    \item Cambiar BPM.
    \item Insertar plugin.
\end{itemize}

%====================================================
\section{Requisitos No Funcionales Avanzados}

\begin{itemize}
    \item Latencia < 50ms.
    \item Arquitectura extensible.
    \item Cobertura mínima 75\%.
    \item Manejo robusto de excepciones.
    \item Compatibilidad Windows.
\end{itemize}

%====================================================
\section{Plan de Pruebas}

\subsection{Unitarias}
Dominio y Application.

\subsection{Integración}
Audio Engine + Plugins.

\subsection{Stress}
Prueba con múltiples pistas activas.

%====================================================
\section{Escalabilidad Futura}

\begin{itemize}
    \item Soporte VST conceptual.
    \item Motor C++ desacoplado.
    \item Render offline avanzado.
    \item Sincronización externa avanzada.
\end{itemize}

%====================================================
\section{Conclusión}

JustBeat-DAW está diseñado como un sistema profesional modular, aplicando principios formales de ingeniería de software a la síntesis digital 8-bit, garantizando extensibilidad, mantenibilidad y evolución futura hacia un producto comercial real.

\end{document}